\subsection{TENSOR PRODUCTS OF g-MODULES}

Let \(\mathrm{E},\mathrm{F}\) be \(\mathfrak{g}\) -modules. For all \(\lambda, \mu \in \mathfrak{h}^*\) , \(\mathrm{E}^\lambda \otimes \mathrm{F}^\mu \subset (\mathrm{E} \otimes \mathrm{F})^{\lambda + \mu}\) (Chap. VII, §1, no. 1, Prop. 2 (ii)). If \(\mathrm{E}\) and \(\mathrm{F}\) are finite dimensional, then \(\mathrm{E} = \sum_{\lambda \in \mathrm{P}} \mathrm{E}^\lambda\) and \(\mathrm{F} = \sum_{\mu \in \mathrm{P}} \mathrm{F}^\mu\) ; consequently,

\[
(\mathrm{E} \otimes \mathrm{F}) ^ {\nu} = \sum_ {\lambda , \mu \in \mathrm{P}, \lambda + \mu = \nu} \mathrm{E} ^ {\lambda} \otimes \mathrm{F} ^ {\mu}.
\]

In other words, equipped with its graduation of type P, \(E \otimes F\) is the graded tensor product of the graded vector spaces E and F.

PROPOSITION 9. Let E,F be finite dimensional simple g-modules, with highest weights \(\lambda,\mu\) , respectively.

\begin{enumerate}
\def\labelenumi{(\roman{enumi})}
\item
  The component of \(\mathrm{E} \otimes \mathrm{F}\) of highest weight \(\lambda + \mu\) is a simple submodule, generated by \((\mathrm{E} \otimes \mathrm{F})^{\lambda + \mu} = \mathrm{E}^{\lambda} \otimes \mathrm{F}^{\mu}\) .
\item
  The highest weight of any simple submodule of \(\mathrm{E} \otimes \mathrm{F}\) is \(\leq \lambda + \mu\) (cf.~§9, Prop. 2).
\end{enumerate}

If \(\alpha, \beta \in \mathrm{P}\) and if \(\mathrm{E}^{\alpha} \otimes \mathrm{F}^{\beta} \neq 0\) , then \(\alpha \leq \lambda\) and \(\beta \leq \mu\) . Consequently, \((\mathrm{E} \otimes \mathrm{F})^{\lambda + \mu}\) is equal to \(\mathrm{E}^{\lambda} \otimes \mathrm{F}^{\mu}\) , and hence is of dimension 1, and \(\lambda + \mu\) is the highest weight of \(\mathrm{E} \otimes \mathrm{F}\) . Every non-zero element of \(\mathrm{E}^{\lambda} \otimes \mathrm{F}^{\mu}\) is primitive. By Prop. 4 of §6, no. 2, the length of the isotypical component of \(\mathrm{E} \otimes \mathrm{F}\) of highest weight \(\lambda + \mu\) is 1.

Remark. Retain the notations of Prop. 9. Let C be the isotypical component of \(E \otimes F\) of highest weight \(\lambda + \mu\) . Then C depends only on E and F and not on the choice of h and the basis B. In other words, let \(h'\) be a splitting Cartan subalgebra of g, \(R'\) the root system of \((\mathfrak{g}, \mathfrak{h}')\) , and \(B'\) a basis of \(R'\) ; let \(\lambda', \mu'\) be the highest weights of E, F relative to \(h'\) and \(B'\) ; let \(C'\) be the isotypical component of \(E \otimes F\) of highest weight \(\lambda' + \mu'\) ; then \(C' = C\) . Indeed, to prove this we can assume, by extension of the base field, that k is algebraically closed. Then there exists \(s \in \operatorname{Aut}_{e}(\mathfrak{g})\) that takes \(h\) to \(h'\) , R to \(R'\) , B to \(B'\) . Let \(S \in \mathbf{SL}(E \otimes F)\) have the properties in Prop. 2 of no. 1. Then \(\mathrm{S}((\mathrm{E} \otimes \mathrm{F})^{\lambda + \mu}) = (\mathrm{E} \otimes \mathrm{F})^{\lambda' + \mu'}\) and \(\mathrm{S}(\mathrm{C}) = \mathrm{C}\) . Hence \((\mathrm{E} \otimes \mathrm{F})^{\lambda' + \mu'} \subset \mathrm{C}' \cap \mathrm{S}(\mathrm{C}) = \mathrm{C}' \cap \mathrm{C}\) , so \(C' = C\) . Thus, to 2 classes of finite dimensional simple g-modules we can associate canonically a third; in other words, we have defined on the set \(S_{g}\) of classes of finite dimensional simple g-modules a composition law. With this structure, \(S_{g}\) is canonically isomorphic to the additive monoid \(P_{++}\) .

COROLLARY 1. Let \((\varpi_{\alpha})_{\alpha \in \mathrm{B}}\) be the family of fundamental weights relative to B. Let \(\lambda = \sum_{\alpha \in \mathrm{B}} m_{\alpha} \varpi_{\alpha} \in \mathrm{P}_{++}\) . For all \(\alpha \in \mathrm{B}\) , let \(\mathrm{E}_{\alpha}\) be a simple \(\mathfrak{g}\) -module of highest weight \(\varpi_{\alpha}\) . In the \(\mathfrak{g}\) -module \(\bigotimes_{\alpha \in \mathrm{B}} (\bigotimes^{m_{\alpha}} \mathrm{E}_{\alpha})\) , the isotypical component of highest weight \(\lambda\) is of length 1.

This follows from Prop. 9 by induction on \(\sum_{\alpha \in \mathrm{B}}m_{\alpha}\) .

COROLLARY 2. Assume that k is R or C or a non-discrete complete ultrametric field. Let G be a Lie group with Lie algebra g. Assume that, for any fundamental representation \(\rho\) of g, there exists an analytic linear representation \(\rho'\) of G such that \(\rho = \mathrm{L}(\rho')\) . Then, for any finite dimensional linear representation \(\pi\) of g, there exists an analytic linear representation \(\pi'\) of G such that \(\pi = \mathrm{L}(\pi')\) .

We use the notations of Cor. 1. There exists a representation \(\sigma\) of G on \(X = \bigotimes_{\alpha \in B} (\bigotimes^{m_{\alpha}} E_{\alpha})\) such that \(L(\sigma)\) corresponds to the \(\mathfrak{g}\) -module structure of X (Chap. III, §3, no. 11, Cor. 3 of Prop. 41). Let C be the isotypical component of X of highest weight \(\lambda\) . In view of Chap. III, §3, no. 11, Prop. 40, it suffices to prove that C is stable under \(\sigma(G)\) . Let \(g \in G\) and \(\varphi = \operatorname{Ad}(g)\) . Then \(\sigma(g)a_X \sigma(g)^{-1} = (\varphi(a))_X\) for all \(a \in \mathfrak{g}\) . On the other hand, \(\varphi\) is an automorphism of \(\mathfrak{g}\) that takes \(\mathfrak{h}\) to \(\mathfrak{h}'\) , R to \(R' = R(\mathfrak{g}, \mathfrak{h}')\) , B to a basis \(B'\) of \(R'\) , and \(\varpi_{\alpha}\) to the highest weight \(\varpi_{\alpha}'\) of \(E_{\alpha}\) relative to \(\mathfrak{h}'\) and \(B'\) (since \(\varphi\) transforms \(E_{\alpha}\) into a \(\mathfrak{g}\) -module isomorphic to \(E_{\alpha}\) ). Hence \(\varphi\) takes \(\lambda\) to \(\sum m_{\alpha} \varpi_{\alpha}'\) . By the Remark above, \(\sigma(g)(C) = C\) .

PROPOSITION 10. Let \(\lambda, \mu \in \mathrm{P}_{++}\) . Let E, F, G be simple \(\mathfrak{g}\) -modules with highest weights \(\lambda, \mu, \lambda + \mu\) . Let \(\mathcal{X}\) (resp. \(\mathcal{X}', \mathcal{X}''\) ) be the set of weights of E (resp. F, G). Then \(\mathcal{X}'' = \mathcal{X} + \mathcal{X}'\) .

We have \(E = \bigoplus_{\nu \in P} E^{\nu}, F = \bigoplus_{\sigma \in P} F^{\sigma}\) , so \(E \otimes F\) is the direct sum of the

\[
(\mathrm{E} \otimes \mathrm{F}) ^ {\tau} = \sum_ {\nu + \sigma = \tau} \mathrm{E} ^ {\nu} \otimes \mathrm{F} ^ {\sigma}.
\]

By Prop. 9, \(\mathbf{G}\) can be identified with a \(\mathfrak{g}\) -submodule of \(\mathrm{E} \otimes \mathrm{F}\) , so \(\mathcal{X}'' \subset \mathcal{X} + \mathcal{X}'\) . We have \(\mathrm{G}^{\tau} = \mathrm{G} \cap (\mathrm{E} \otimes \mathrm{F})^{\tau}\) , and it is enough to show that, for \(\nu \in \mathcal{X}\) and \(\sigma \in \mathcal{X}'\) , we have \(\mathrm{G} \cap (\mathrm{E} \otimes \mathrm{F})^{\nu + \sigma} \neq 0\) . Let \((e_1, \ldots, e_n)\) (resp. \((f_1, \ldots, f_p))\) be a basis of \(\mathrm{E}\) (resp. \(\mathrm{F}\) ) consisting of elements each of which belong to some \(\mathrm{E}^\nu\) (resp. \(\mathrm{F}^\sigma\) ), and such that \(e_1 \in \mathrm{E}^\lambda\) (resp. \(f_1 \in \mathrm{F}^\mu\) ). The \(e_i \otimes f_j\) form a basis of \(\mathrm{E} \otimes \mathrm{F}\) . Suppose that the result to be proved is false. Then there exists a pair \((i,j)\) such that the coordinate of index \((i,j)\) of every element of \(\mathrm{G}\) is zero. Let \(\mathrm{U}\) be the enveloping algebra of \(\mathfrak{g}\) , \(\mathrm{U}'\) the dual of \(\mathrm{U}\) , \(c\) the coproduct of \(\mathrm{U}\) . For all \(u \in \mathrm{U}\) , let \(x_i(u)\) (resp. \(y_j(u)\) ) be the coordinate of \(u(e_1)\) (resp. \(u(f_1)\) ) of index \(i\) (resp. \(j\) ); let \(z_{ij}(u)\) be the coordinate of index \((i,j)\) of \(u(e_1 \otimes f_1)\) . Then \(x_i, y_j, z_{ij} \in \mathrm{U}'\) . Now \(e_1\) generates the \(\mathfrak{g}\) -module \(\mathrm{E}\) , so \(x_i \neq 0\) , and similarly \(y_j \neq 0\) . By the definition of the \(\mathfrak{g}\) -module \(\mathrm{E} \otimes \mathrm{F}\) (Chap. I, §3, no. 2), if \(c(u) = \sum u_s \otimes u_s'\) , we have

\[
z _ {i j} (u) = \sum_ {s} x _ {i} (u _ {s}). y _ {j} (u _ {s} ^ {\prime}) = \langle c (u), x _ {i} \otimes y _ {j} \rangle .
\]

In other words, \(z_{ij}\) is the product of \(x_{i}\) and \(y_{j}\) in the algebra \(U'\) . But this algebra is an integral domain (Chap. II, §1, no. 5, Prop. 10), so \(z_{ij} \neq 0\) . Since \(u(e_{1} \otimes f_{1}) \in \mathrm{G}\) for all \(u \in U\) , this is a contradiction.

